
The Perseus multiple systems were observed with ALMA during
Cycle 5 with 5 executions throughout 2017 December, 2018 January, and 2018 September with 43-45 
antennas operating and sampling baselines between $\sim$15 - $\sim$2500 meters. 
The observations were executed within $\sim$1.1 hour block
and the total time spent on each source was $\sim$17.6-20.2 minutes. 
The precipitable water vapor ranged from 0.54 - 2.57~mm during the different executions.
The phase calibrator was J0336+3218 in all observations and the bandpass and flux calibrator was either
J0238+1636 or J0237+2848.
The correlator was configured for one baseband to observe a 2 GHz continuum band 
centered at 232.5 GHz and observed in TDM mode with 128 channels. Then the second baseband
was configured with two 234~MHz spectral windows with 488 kHz channels centered on $^{12}$CO ($J=2\rightarrow1$) and
N$_2$D$^{+}$ ($J=3\rightarrow2$) and $^{13}$CS ($J=5\rightarrow4$). The third baseband was configured with four
59~MHz spectral windows with 122 kHz channels centered on H$_2$CO ($J=3_{03}\rightarrow2_{02}$), ($J=3_{22}\rightarrow2_{21}$), 
($J=3_{21}\rightarrow2_{20}$), and SO ($J_N = 6_5\rightarrow5_4$). The fourth baseband was configured with two 59 MHz spectral windows with 61 kHz channels
centered on $^{13}$CO ($J=2\rightarrow1$) and C$^{18}$O ($J=2\rightarrow1$). However, we will only discuss the $^{12}$CO ($J=2\rightarrow1$)
data in this paper.

The raw visibility data were reduced by the ALMA pipeline integrated within CASA 5.4.0. We additionally performed
self-calibration on the continuum data to increase the signal to noise ratio. We performed 
self-calibration of the continuum data on all sources also using CASA 5.6.1, using two rounds of phase-only self-calibration where the first solution
interval corresponded to the length of a single scan and the second solution interval corresponded 
to 12.1s or 6.05s; 6.05s corresponds to the length of a single integration. Following phase-only
self-calibration, we also applied amplitude self-calibration, using a solution interval that was 
the length of a scan, and we used the option \textt{solnorm=True} to normalize the amplitude solutions to 1, rather than
allowing them to scale freely, which prevents the flux density from being significantly altered erroneously.
Following the completion of self-calibration on the continuum, the solutions were also applied to the 
spectral line spectral windows.

The data were imaged using a custom recipe of the ALMA imaging pipeline in CASA version 5.4.0, using Briggs weighting with robust=0.5 for all images. 
The continuum images were produced using standard 
calls to \texttt{hif\_makeimlist} and \texttt{hif\_makeimages}. We used options of the \texttt{hif\_editimlist}
task to image the spectral windows with the desired channel binning and frequency range. For $^{12}$CO, we used 0.66~\kms\ channels
and imaged $\pm$100~\kms\ of the source velocity. The $^{12}$CO cubes typically had 800$\times$800 pixels, a pixel scale of 0\farcs053, and a synthesized beam of 0\farcs38$\times$0.\farcs30. The continuum images were typically 1120$\times$1120 pixels and a pixel scale of 0\farcs04 and the synthesized beam was 0\farcs28$\times$0\farcs21. The noise in the 1.3~mm continuum was $\sim$0.2~mJy~beam$^{-1}$ and the noise in the $^{12}$CO cubes was 3.6~mJy~beam$^{-1}$.

One source, L1448 IRS3C, was not self-calibrated in the same way as the others because it had been unintentionally left out of the orginal processing. Instead
of manually processing the data, we imaged those data with the ALMA pipeline included in CASA 6.6.1 using automated self-calibration through
the imaging service available on the NRAO archive. Comparison of self-calibration done through the automated imaging
service versus the manual self-calibration showed very similar results, and the ALMA imaging pipeline was executed in a very similar way to all the other images as well.


